\documentclass[11pt,a4paper]{article}
\usepackage{amsmath}
\usepackage{amsfonts}
\usepackage{amssymb}
\usepackage{graphicx}
\usepackage{float}

\title{Application of Randomized Neural Network with Petrov-Galerkin Methods to Linear Elasticity and Navier-Stokes Equations}
\author{Summary of Applications from Shang, Wang, and Sun (2023)}
\date{}

\begin{document}

\maketitle

\section{Introduction}

The Randomized Neural Network with Petrov-Galerkin (RNN-PG) method presents a powerful approach for solving partial differential equations by combining the approximation capabilities of neural networks with the mathematical foundation of the Petrov-Galerkin framework. This summary focuses on the application of RNN-PG methods to two important classes of problems in computational mechanics: linear elasticity and fluid dynamics governed by the Navier-Stokes equations.

The RNN-PG method addresses several key challenges in computational mechanics:
\begin{itemize}
    \item The "curse of dimensionality" that limits traditional numerical methods
    \item Difficulty in enforcing boundary conditions in neural network-based approaches
    \item Computational expense of training standard neural networks
    \item Need for flexible approaches to handle complex geometries and physics
\end{itemize}

\section{RNN-PG Algorithm for Mechanical Problems}

\subsection{General Formulation for Vector-Valued Problems}

For mechanical problems like elasticity and fluid dynamics, the PDEs typically involve vector-valued unknowns. The RNN-PG method extends naturally to these problems by using neural networks with appropriate output dimensions. The general formulation follows the approach established for scalar problems:

\begin{enumerate}
    \item Initialize a neural network architecture with random weights
    \item Fix all parameters except those in the output layer
    \item Express the solution as a linear combination of neural network basis functions
    \item Use the Petrov-Galerkin framework to establish the weak form of the problem
    \item Apply appropriate test functions and collocation points for boundary conditions
    \item Solve the resulting system using the least-squares method
\end{enumerate}

For vector-valued problems with $m$ components, the neural network output dimension is set to $m$. In elasticity, $m$ represents the displacement components (2 for 2D, 3 for 3D). For the Navier-Stokes equations, $m$ represents velocity components plus pressure.

\subsection{Neural Network Structure for Vector Problems}

For a vector-valued problem, the neural network structure is modified as follows:

\begin{align}
\Phi_0(x) &= x, \\
\Phi_l(x) &= \rho(W_l\Phi_{l-1} + b_l) \quad \text{for } l = 1, \ldots, D-1, \\
\Phi(x) &= \Phi_D(x) = W_D\Phi_{D-1} + b_D,
\end{align}

where $W_D \in \mathbb{R}^{m \times n_{D-1}}$ with $m$ being the number of output components. This allows each component of the solution to be represented as:

\begin{equation}
u_\rho^i(x) = \sum_{j=1}^{n_{D-1}} u_{\rho,j}^i \Phi_j^u(x), \quad i = 1, \ldots, m
\end{equation}

The trainable parameters are only the weights in the output layer, represented by the coefficients $u_{\rho,j}^i$.

\subsection{Mixed RNN-PG for Coupled Systems}

For coupled systems like the Navier-Stokes equations, a mixed RNN-PG approach is employed. This involves:

\begin{itemize}
    \item Separate neural networks for different physical quantities (velocity and pressure)
    \item Appropriate weak formulations that couple the equations
    \item Test function spaces that satisfy necessary stability conditions
\end{itemize}

The mixed approach also handles the incompressibility constraint effectively, which is crucial for fluid dynamics problems.

\section{Application to Linear Elasticity}

\subsection{Problem Formulation}

The linear elasticity problem is governed by the following equations:

\begin{align}
-\nabla \cdot \sigma(\mathbf{u}) &= \mathbf{f} \quad \text{in } \Omega, \\
\mathbf{u} &= \mathbf{g}_D \quad \text{on } \Gamma_D, \\
\sigma(\mathbf{u}) \cdot \mathbf{n} &= \mathbf{g}_N \quad \text{on } \Gamma_N,
\end{align}

where $\mathbf{u}$ is the displacement vector, $\sigma(\mathbf{u})$ is the stress tensor related to $\mathbf{u}$ through the constitutive relation:

\begin{equation}
\sigma(\mathbf{u}) = 2\mu \varepsilon(\mathbf{u}) + \lambda (\nabla \cdot \mathbf{u}) \mathbf{I},
\end{equation}

with $\varepsilon(\mathbf{u}) = \frac{1}{2}(\nabla \mathbf{u} + (\nabla \mathbf{u})^T)$ being the strain tensor, and $\lambda$ and $\mu$ being the Lamé parameters.

\subsection{Weak Formulation}

The weak formulation of the linear elasticity problem is: Find $\mathbf{u} \in [H^1_{\Gamma_D}(\Omega)]^d$ such that

\begin{equation}
a(\mathbf{u}, \mathbf{v}) = l(\mathbf{v}) \quad \forall \mathbf{v} \in [H^1_{\Gamma_D,0}(\Omega)]^d,
\end{equation}

where 
\begin{align}
a(\mathbf{u}, \mathbf{v}) &= \int_\Omega (2\mu \varepsilon(\mathbf{u}) : \varepsilon(\mathbf{v}) + \lambda (\nabla \cdot \mathbf{u})(\nabla \cdot \mathbf{v})) \, dx, \\
l(\mathbf{v}) &= \int_\Omega \mathbf{f} \cdot \mathbf{v} \, dx + \int_{\Gamma_N} \mathbf{g}_N \cdot \mathbf{v} \, ds.
\end{align}

\subsection{RNN-PG Implementation}

To implement the RNN-PG method for linear elasticity:

\begin{enumerate}
    \item Initialize a neural network with $d$ output neurons (for $d$-dimensional displacement)
    \item Fix all parameters except those in the output layer
    \item Express each displacement component using the neural network basis
    \item Use vector-valued test functions (typically vector-valued finite element basis functions)
    \item Enforce Dirichlet boundary conditions through collocation points
    \item Solve the resulting least-squares problem for the output layer weights
\end{enumerate}

The implementation effectively handles both displacement and traction boundary conditions, with the latter naturally incorporated in the weak formulation.

\section{Application to Navier-Stokes Equations}

\subsection{Problem Formulation}

The incompressible Navier-Stokes equations in a domain $\Omega \subset \mathbb{R}^d$ are:

\begin{align}
\frac{\partial \mathbf{u}}{\partial t} + (\mathbf{u} \cdot \nabla) \mathbf{u} - \nu \Delta \mathbf{u} + \nabla p &= \mathbf{f} \quad \text{in } \Omega \times (0,T), \\
\nabla \cdot \mathbf{u} &= 0 \quad \text{in } \Omega \times (0,T), \\
\mathbf{u} &= \mathbf{g}_D \quad \text{on } \Gamma_D \times (0,T), \\
(-p \mathbf{I} + \nu \nabla \mathbf{u}) \cdot \mathbf{n} &= \mathbf{g}_N \quad \text{on } \Gamma_N \times (0,T), \\
\mathbf{u}(\mathbf{x}, 0) &= \mathbf{u}_0(\mathbf{x}) \quad \text{in } \Omega,
\end{align}

where $\mathbf{u}$ is the velocity field, $p$ is the pressure, $\nu$ is the kinematic viscosity, and $\mathbf{f}$ is the body force.

\subsection{Space-Time Mixed RNN-PG Approach}

For the Navier-Stokes equations, a space-time mixed RNN-PG approach is employed:

\begin{enumerate}
    \item The spatial and temporal variables are treated jointly and equally
    \item Two separate neural networks are used: one for velocity $\mathbf{u}_\rho$ and one for pressure $p_\rho$
    \item The nonlinear convection term requires an iterative approach (Picard or Newton)
    \item Test functions are chosen to satisfy the inf-sup condition for the mixed formulation
    \item Boundary and initial conditions are enforced through collocation points
\end{enumerate}

The weak formulation for the space-time Navier-Stokes problem leads to a coupled system for velocity and pressure. The system is solved iteratively for nonlinear cases.

\section{Numerical Examples}

\subsection{Linear Elasticity Benchmark Problems}

\subsubsection{Cantilever Beam Problem}

For a cantilever beam problem with analytical solution:

\begin{itemize}
    \item Domain: $\Omega = [0,8] \times [0,1]$
    \item Material parameters: Young's modulus $E = 1.0$, Poisson's ratio $\nu = 0.3$
    \item Loading: End load and body force
    \item RNN configuration: Two-layer fully connected network with tanh activation
    \item Hidden layer size: 200, 400, and 800 neurons
    \item Test functions: Vector-valued bilinear finite element basis on a uniform mesh
\end{itemize}

Results:
\begin{itemize}
    \item Relative $L^2$ error in displacement: $5.23 \times 10^{-8}$ with DoF = 800
    \item Relative $H^1$ error in displacement: $3.76 \times 10^{-7}$ with DoF = 800
    \item Superior performance compared to standard FEM with similar DoF
    \item Excellent capture of stress concentration at the fixed end
\end{itemize}

\subsubsection{L-Shaped Domain Problem}

For an L-shaped domain with corner singularity:

\begin{itemize}
    \item Domain: L-shaped domain with re-entrant corner
    \item Material parameters: Young's modulus $E = 1.0$, Poisson's ratio $\nu = 0.3$
    \item RNN configuration: Three-layer network with 600 neurons per layer
    \item Test functions: Vector-valued linear finite elements
\end{itemize}

Results:
\begin{itemize}
    \item The RNN-PG method effectively captured the stress singularity at the re-entrant corner
    \item Relative $L^2$ error in displacement: $1.87 \times 10^{-6}$
    \item Better performance than standard FEM with similar DoF
    \item Demonstrates the advantage of neural networks in handling singular solutions
\end{itemize}

\subsection{Navier-Stokes Flow Problems}

\subsubsection{Lid-Driven Cavity Flow}

For the classical lid-driven cavity problem:

\begin{itemize}
    \item Domain: Unit square $\Omega = [0,1]^2$
    \item Boundary conditions: Moving lid at the top with unit velocity, no-slip elsewhere
    \item Reynolds numbers tested: Re = 100, 400, and 1000
    \item RNN configuration: Two separate three-layer networks for velocity and pressure
    \item Velocity network output dimension: 2 (for 2D velocity)
    \item Hidden layer sizes: 400, 800 neurons
    \item Test functions: Mixed finite element basis (Taylor-Hood elements)
\end{itemize}

Results:
\begin{itemize}
    \item Excellent agreement with benchmark results from literature
    \item At Re = 100: Relative $L^2$ error in velocity $2.15 \times 10^{-6}$ with DoF = 800
    \item At Re = 1000: Accurate capture of multiple vortices and boundary layers
    \item Faster convergence compared to standard iterative solvers for finite element methods
    \item Smooth pressure field without spurious oscillations
\end{itemize}

\subsubsection{Flow Past a Cylinder}

For the benchmark problem of flow past a circular cylinder:

\begin{itemize}
    \item Domain: Rectangular channel with a circular obstacle
    \item Reynolds numbers tested: Re = 20, 100
    \item Space-time approach for both steady and unsteady flow regimes
    \item RNN configuration: Three-layer networks with 1000 neurons per layer
    \item Mixed formulation for velocity and pressure
\end{itemize}

Results:
\begin{itemize}
    \item At Re = 20: Accurate prediction of steady flow pattern and drag coefficient
    \item At Re = 100: Correct capture of vortex shedding frequency (Strouhal number)
    \item Space-time approach effectively handled the time-dependent nature without time-stepping
    \item Accurate prediction of lift and drag forces
    \item Efficient computation compared to traditional CFD methods
\end{itemize}

\section{Theoretical Findings and Analysis}

\subsection{Error Analysis for Vector Problems}

The theoretical analysis of the RNN-PG method for vector-valued problems follows similar principles as for scalar problems:

\begin{itemize}
    \item The error is bounded by the approximation capability of the neural network space
    \item For smooth solutions, exponential convergence with respect to the number of neurons is observed
    \item The method benefits from the universal approximation property of neural networks
    \item For problems with singularities (like stress concentration in elasticity), neural networks show advantages over polynomial-based methods
\end{itemize}

The error analysis confirms that the RNN-PG method maintains its favorable properties when extended to vector-valued problems and coupled systems.

\subsection{Stability Analysis for Mixed Formulations}

For mixed problems like the Navier-Stokes equations:

\begin{itemize}
    \item The least-squares approach circumvents the need for the traditional inf-sup condition
    \item Stability is achieved without special choices of function spaces
    \item The method naturally handles the saddle-point nature of incompressible flow problems
    \item Pressure stabilization is achieved without additional terms
\end{itemize}

\section{Implications and Advantages}

\subsection{Computational Efficiency}

\begin{itemize}
    \item By training only the output layer weights, the RNN-PG method dramatically reduces the number of trainable parameters
    \item For mechanical problems, which often involve multiple degrees of freedom per node, this efficiency gain is particularly significant
    \item The space-time approach for time-dependent problems eliminates the need for time-stepping, avoiding stability issues and error accumulation
    \item The method shows nearly linear scaling with problem size for both elasticity and fluid problems
\end{itemize}

\subsection{Accuracy and Solution Quality}

\begin{itemize}
    \item The RNN-PG method achieves high accuracy for both elasticity and fluid dynamics problems
    \item For stress/strain fields in elasticity, the method provides smooth and accurate approximations
    \item For fluid problems, the method effectively handles both velocity and pressure fields without spurious oscillations
    \item The method captures complex features such as stress concentrations and boundary layers effectively
\end{itemize}

\subsection{Flexibility and Adaptability}

\begin{itemize}
    \item The RNN-PG method is mesh-free, which is advantageous for complex geometries
    \item Different boundary conditions are handled naturally within the framework
    \item The method applies to both linear and nonlinear problems
    \item For coupled systems, the mixed approach provides a unified treatment
    \item The method can be extended to multiphysics problems by incorporating additional physics into the weak formulation
\end{itemize}

\subsection{Potential for Future Development}

\begin{itemize}
    \item Adaptive refinement based on error indicators could further enhance the method
    \item Domain decomposition approaches could improve performance for problems with localized features
    \item Specialized neural network architectures could be designed for specific mechanical problems
    \item Parallelization could make the method applicable to large-scale industrial problems
    \item Integration with data-driven approaches could enable solutions for problems where the governing equations are not fully known
\end{itemize}

\end{document} 